% Cover letter for J. R. Soc. Interface; same text as cover_letter.md (keep the two in step).
% Bracketed bold items are for the author to decide or write before submission.
% Build: latexmk -pdf cover_letter.tex
\documentclass[11pt]{article}
\usepackage[a4paper,margin=2.5cm]{geometry}
\usepackage[T1]{fontenc}
\usepackage{lmodern}
\usepackage{url}
\setlength{\parindent}{0pt}
\setlength{\parskip}{0.6em}
\newcommand{\byauthor}[1]{\textbf{[#1]}}
\pagestyle{empty}
\begin{document}

Dear Editor,

Please consider our manuscript ``Silent but deadly? Physics says seldom, memory says mostly'' for publication as a
Research Article in \emph{Journal of the Royal Society Interface}.

The belief that a silent release of flatus smells worse than an audible one is recorded in many languages, yet no
mechanistic test of it exists. The closest prior work is a conference abstract on the physics of flatulence (Ancalle
et al.\ 2021). We follow one passage through a single chain: chemistry, sound, dispersion, perception and memory.
Three results reach beyond the topic.

\begin{enumerate}
  \item \textbf{Routing by the human convective boundary layer.} In simulations of a sofa, a standing pair and a closed
  lift, the warm air rising along a body takes a slow release next to it and lets a fast one pass, so the momentum of
  the release, relative to that of the boundary layer, decides which nearby nose it reaches. The same set-up
  reproduces measured breathing-zone intake from a source at a seated manikin's groin within a factor of 1.5. This
  bears on any near-body source, such as respiratory, skin and clothing emissions in shared indoor spaces.
  \item \textbf{A peak bound for linear transport.} We prove that a slower release of the same substance cannot give a
  higher peak concentration, in any room or flow. Only the non-linear near field lets it win, by a factor bounded by a
  velocity ratio, which we call the Silencer number.
  \item \textbf{Selection across two senses.} With passages as emitted and settings weighted by how often people share
  them, a bystander finds the silent passage worse in 0.5\% of pairs and the audible one in 6.9\%; only in a closed
  lift, for the person behind the releaser, are the two about equal. Yet because a silent passage is noticed only if
  it is smelled, a remembered silent passage seems the worse one in 88\% of comparisons. When an event can be noticed
  through either of two senses, selection on noticing alone can create a false sensory belief. The explanation makes
  testable predictions: releasers should hold the belief less firmly than bystanders, and people without a sense of
  smell should never remember a silent passage.
\end{enumerate}

The work joins fluid dynamics, chemistry, acoustics, sensory science and cognition. We think it suits
\emph{Interface} because each part can be read on its own, and because the question shows how a physical and a
cognitive explanation can be separated. The headline is checked against the scene weights, the observer models,
large-eddy peak factors, mesh refinement and clothing; its direction holds in every case but one adverse pairing of
peak factors, and the remembered share stays above 80\% in all of them.

Equations and proofs are in appendices, and technical details are in the electronic supplementary material. Code and
processed data are openly available at \url{https://github.com/ccy5123/sbd}; an archived release with a DOI (Zenodo)
will accompany the revised manuscript. \byauthor{A preprint is posted on engrXiv (DOI \dots).} The results have not been
reported in the media. The manuscript is not under consideration elsewhere. \byauthor{AI-use statement, written by the
author.}

Yours sincerely,\\[0.4em]
Chan Young Joe\\
Independent Researcher, Incheon, Republic of Korea\\
ORCID 0009-0007-5822-6714

\end{document}
